\section{A total-positivity boundary and two actual character constructions}

\begin{theorem}[NN-A007: the Newman threshold is a PF boundary]
Define the entire function
\[
 F_t(w)=\frac{H_t(i\sqrt w)}{H_t(0)}=\sum_{k\ge0}c_k(t)w^k,
 \qquad
 c_k(t)=\frac{\int_0^\infty e^{tu^2}\Phi(u)u^{2k}\,du}
 {(2k)!H_t(0)}.
\]
The square-root notation denotes the entire even-power series, so it
introduces no branch choice. Every $c_k(t)$ is strictly positive for
real $t$. Under A, the infinite Toeplitz matrix
\[
 T_t=(c_{j-i}(t))_{i,j\ge0},\qquad c_k=0\quad(k<0),
\]
is totally nonnegative if and only if $t\ge0$. Thus for each $t<0$
some finite minor is negative, although every entry on or above its
diagonal is positive. At each $t\ge0$,
\[
 F_t(w)=\prod_j(1+w/x_j(t)^2),
\]
where $x_j(t)>0$ lists the positive zeros of $H_t$, with multiplicities.
\end{theorem}
\begin{proof}
Positivity and convergence follow directly from the theta integral.
The order estimate $\log\max_{|z|\le R}|H_t(z)|=O_t(R\log R)$
in workbench22, \texttt{sec:rh-newman-collision}, implies that $F_t$
has order at most $1/2$. When all zeros of $H_t$ are real, pairing
$x_j$ with $-x_j$ gives the displayed genus-zero product; the possible
zero-free exponential factor is constant because the order of $F_t$
is less than one, and $F_t(0)=1$ fixes it. Moreover
$\sum_jx_j^{-2}<\infty$.

Each factor $1+bw$, $b\ge0$, has upper-bidiagonal Toeplitz matrix
$I+bS$, which is totally nonnegative: every nonzero minor has its
order-preserving matching and nonnegative product. Multiplication of
generating series is multiplication of their Toeplitz matrices.
Cauchy--Binet therefore gives nonnegative minors for every finite
product. Local uniform convergence gives coefficient convergence and
hence convergence of each finite determinant, proving the forward
direction without invoking a classification theorem.

For the converse use the known Edrei theorem, in the precise original
form of Proposition A on pp.86--87 and its proof in Section 3, p.92,
of Edrei, \emph{Proof of a conjecture of Schoenberg}, Canad. J. Math.
5 (1953), 86--94, DOI 10.4153/CJM-1953-010-3: an entire generating
function of a totally nonnegative sequence has all its zeros on the
negative real axis. A zero $z$ of $H_t$ gives a zero $w=-z^2$ of
$F_t$. If all such $w$ are negative real, every $z$ is real. Thus PF
is equivalent to all-real zeros. The threshold property and A together
with Rodgers--Tao give exactly $t\ge0$. Failure of total nonnegativity
means, by its definition, failure at a finite minor; coefficients are
real, so that minor is strictly negative.
\end{proof}

The distinction between positive coefficients and all minors is material:
entrywise positivity persists throughout the nonreal-zero phase. The
negative certificate has finite size for each fixed negative time, but
this argument supplies neither its size nor an algorithmic rate as
$t\uparrow0$. It therefore gives no complexity upper bound.

\begin{theorem}[NN-A008: zeta parameters in symmetric, unitary, and flag geometry]
Under A, let $\gamma_j>0$ list the positive zeta ordinates with
multiplicities and put
\[
 C=\sum_j\gamma_j^{-2},\qquad
 \beta_j=\gamma_j^{-2}/C,\qquad
 E(w)=\prod_j(1+\beta_jw)=\sum_{k\ge0}e_k(\beta)w^k.
\]
Then $C$ is finite and positive, every $0<\beta_j<1$, and
$\sum\beta_j=1$. The following constructions are well defined.
\begin{enumerate}
\item A normalized central positive-definite character of $S_\infty$ is
\[
 \chi_S(\sigma)=\prod_{k\ge2}
 \left((-1)^{k-1}\sum_j\beta_j^k\right)^{m_k(\sigma)}.
\]
It has Thoma parameters $(\alpha,\beta,\gamma)=(0,\beta,0)$.
\item A normalized continuous central positive-definite character on
$U(\infty)=\bigcup_nU(n)$ is
\[
 \chi_U(U)=\prod_j\det((1-\beta_j)I_n+\beta_jU),\qquad U\in U(n).
\]
These are two different characters, not restriction of a single formula.
On a permutation cycle of length $k$, the latter has factor
\[
 \prod_j\{(1-\beta_j)^k-(-\beta_j)^k\},
\]
whereas the former has factor $(-1)^{k-1}\sum_j\beta_j^k$.
\item Every finite upper-triangular Toeplitz truncation
$T^{(n+1)}=(e_{j-i}(\beta))_{0\le i,j\le n}$ has strictly positive
minors whenever upper triangularity does not force them to vanish.
For $1\le i\le n$, put
\[
 D_i=\det T^{(n+1)}_{\{0,\ldots,i-1\},\{n+1-i,\ldots,n\}},
 \quad D_0=D_{n+1}=1,\quad
 q_i=\frac{D_{i-1}D_{i+1}}{D_i^2}>0.
\]
These are literal positive coordinates in the finite Toeplitz
parametrization associated with quantum cohomology of the complete flag
variety $\mathrm{Fl}_{n+1}$, by the specified Rietsch theorem below.
\end{enumerate}
\end{theorem}
\begin{proof}
The paired Hadamard product for $\Xi(z)=\xi(1/2+iz)$ and the zero-count
bound give $\sum\gamma_j^{-2}<\infty$; its infinitely many positive
terms give $0<\beta_j<1$.
For (1), let $D$ be the positive trace-class diagonal operator with
eigenvalues $\beta_j$ on $\ell^2$, and let $U_\sigma$ permute the
factors of $(\ell^2)^{\otimes n}$. Direct summation in the tensor basis
gives
\[
 \operatorname{Tr}(D^{\otimes n}U_\sigma)
 =\prod_{k\ge1}(\operatorname{Tr}D^k)^{m_k(\sigma)}.
\]
Multiplication by $\operatorname{sgn}\sigma$ gives exactly the stated
formula, since $\operatorname{Tr}D=1$. It is compatible with adjoining
fixed points. For any finite complex coefficients $a_i$ and permutations
$\sigma_i$, its positive-definiteness is
\[
 \sum_{i,j}\overline{a_i}a_j\chi_S(\sigma_i^{-1}\sigma_j)
 =\operatorname{Tr}(D^{\otimes n}A^*A)\ge0,
 \quad A=\sum_j a_j\operatorname{sgn}(\sigma_j)U_{\sigma_j}.
\]
The tensor density commutes with all permutations, so the character is
central. The displayed parameters coincide with the standard Thoma
formula; its extremality is the known classification theorem, with exact
modern primary locator Okounkov, arXiv:math/9803037, Section 2.3,
Theorem (Thoma), printed p.23. No classification is needed for the
explicit existence proof just given.

For (2), at a fixed $0\le b\le1$ the determinant expands as
\[
 \det((1-b)I_n+bU)
 =\sum_{k=0}^n b^k(1-b)^{n-k}\operatorname{Tr}(\bigwedge^kU).
\]
It is a nonnegative combination of representation characters, normalized
at the identity because the dimensions sum to $(b+1-b)^n=1$.
Products remain positive definite. For fixed $n$, uniform convergence
follows from $|\det(I_n+\beta_j(U-I_n))-1|\le C_n\beta_j$ and
$\sum\beta_j<\infty$; the finitely many early factors cause no issue,
including any zeros. Hence the limit is continuous and positive
definite. The embedding $U\mapsto U\oplus1$ leaves every factor
unchanged. For a permutation $k$-cycle the eigenvalues are the $k$th
roots of unity, whose product gives $(1-b)^k-(-b)^k$.

For (3), use the finite products of the bidiagonal factors in the
preceding proof. A minor with increasing row indices $r_i$ and column
indices $c_i$ is forced zero unless $r_i\le c_i$ for all $i$.
When these inequalities hold, realize a positive term in the iterated
Cauchy--Binet formula by moving the rightmost selected index from $r_d$
to $c_d$, one unit at a time, then the next rightmost to its target,
and so on. At every step the selected indices are distinct and the
corresponding bidiagonal minor is either one or a positive $\beta_j$.
After finitely many steps all targets are reached. Further factors
have a diagonal term equal to one, so every longer product retains a
fixed strictly positive lower bound for this minor. Passing to the
limit proves strictness. The $D_i$ therefore are positive.
Rietsch, \emph{Totally Positive Toeplitz Matrices: Classical and Modern},
ICM2026, DOI 10.1137/25M1806727, Theorem 2.1 and equations (2.1)--(2.3),
identifies the positive Toeplitz space with the positive diagonal torus
using $d_i=D_i/D_{i-1}$ and $d_{n+1}=1/D_n$. Its successive-ratio
coordinates are exactly the displayed $q_i$. Section 2.1 identifies
these with the quantum parameters for $\mathrm{Fl}_{n+1}$.
\end{proof}

For a fully explicit finite representation measure, let
$\rho_n(\sigma)=\operatorname{sgn}(\sigma)U_\sigma$ and let
$P_\lambda$ be its central orthogonal projection onto the $S_n$
isotypic component of type $\lambda$. Then
\[
 p_n(\lambda)=\operatorname{Tr}(D^{\otimes n}P_\lambda)\ge0,
 \qquad\sum_{\lambda\vdash n}p_n(\lambda)=1,
\]
and $\chi_S|_{S_n}$ is the mixture of normalized irreducible
characters with these weights. Orthogonal-projection positivity and
$\sum P_\lambda=I$ prove both assertions. The standard Frobenius formula
identifies these weights as $f^\lambda s_{\lambda'}(\beta)$; the
factor $f^\lambda$ is essential when using normalized irreducible
characters.

This gives actual algebraic-group flag varieties and representation
objects with explicit real parameters. The parameters are infinite
analytic zero data; they are not nonnegative integer structure constants,
not an asserted canonical basis for GCT, and not Frobenius eigenvalues
of a finite Weil object. The finite quantum-cohomology dictionary is a
specified imported parametrization; it does not create an arithmetic
model of $\operatorname{Spec}\mathbb Z$ or prove a GCT positivity
hypothesis. In particular it is compatible with the supplied GCT
counterexamples.
